\documentclass[11pt]{article}

\usepackage[a4paper,margin=1in]{geometry}
\usepackage{amsmath,amssymb,amsthm}
\usepackage{booktabs}
\usepackage{enumitem}
\usepackage{hyperref}

\title{The Consistency/Coherence Square\\
\large Working note --- intuition pump, not a theorem}
\author{Patrick S.\ Mahon \\ \small Development log entry, 2026-05-15}
\date{}

\begin{document}
\maketitle

\section*{Two faces of consistency}

Consistency has a syntactic and a semantic face:
\begin{itemize}[nosep]
  \item \textbf{Syntactic:} No derivation of $\bot$ from the axioms.
  \item \textbf{Semantic:} There exists a model.
\end{itemize}
The completeness theorem connects them: for first-order logic, these coincide.

\medskip\noindent
\textbf{Question.} Does coherence have two faces?  Or is it inherently
semantic --- the failure mode being non-compact, hence inaccessible to
any finitary proof calculus?

\section*{The punnet square}

Rather than seeking a syntactic face of coherence directly, consider
the $2 \times 2$ of consistent/inconsistent $\times$ coherent/incoherent:

\bigskip
\begin{center}
\begin{tabular}{l|p{5.5cm}|p{5.5cm}|}
\cline{2-3}
& \textbf{Consistent} & \textbf{Inconsistent} \\
\hline
\multicolumn{1}{|l|}{\textbf{Coherent}}
& Classical probability.  Local data embeds in one Boolean algebra;
  CE holds; $\sigma$-additive measure exists.
& Quantum probability.  Incompatible observables (no single Boolean
  embedding), but Gleason gives a global $\sigma$-additive state. \\
\hline
\multicolumn{1}{|l|}{\textbf{Incoherent}}
& The ultrafilter charge.  No contradiction --- every first-order test
  passes --- but $\sigma$-additivity fails.  CE fails.
& Pt\'{a}k's ``exotic logics'' (1987).  Non-Boolean OML admitting
  finitely additive states but no $\sigma$-additive state.
  Obstruction can come from the center $C(L)$. \\
\hline
\end{tabular}
\end{center}

\section*{Candidate axes}

\begin{itemize}[nosep]
  \item \textbf{Consistency axis} $=$ distributivity / Boolean
    embeddability.  ``Can the local data coexist in a single
    classical context?''
  \item \textbf{Coherence axis} $=$ state extension /
    $\sigma$-additivity / CE-type admissibility.  ``Does the local
    data force a global realization of the target property?''
\end{itemize}

\section*{Programme mapping}

The programme's papers correspond to slices of the square:
\begin{itemize}[nosep]
  \item \textbf{Paper I:} Top row (Boolean).  Distinguishes coherent
    from incoherent by identifying CE as the admissibility condition.
  \item \textbf{Paper II:} Left column (state exists).  Distinguishes
    consistent from inconsistent --- i.e., Boolean (EA/PR/VDR all
    available) from non-Boolean (VDR blocked by Kochen--Specker).
  \item \textbf{The diagonal} from top-left to bottom-right traces the
    full filtration: $\mathrm{VDR} \to \mathrm{PR} \to \mathrm{EA}
    \to \text{nothing}$.
\end{itemize}

\section*{What's excluded from this square}

\begin{itemize}[nosep]
  \item \textbf{Kochen--Specker:} about dispersion-free states (VDR),
    not $\sigma$-additive states.  A different global property $P$.
    KS is a compact/finitary obstruction (finite colouring argument),
    closer to inconsistency than to incoherence.
  \item \textbf{Abramsky--Brandenburger contextuality:} about joint
    distributions on compatible sets, closer to VDR level.
  \item \textbf{De Finetti coherence:} the EA level (finite
    additivity) --- below this square.
\end{itemize}

\section*{The key tension}

The horizontal axis is labeled ``Boolean / non-Boolean'' not
``consistent / inconsistent.''  Whether non-Booleanness constitutes
a species of ``inconsistency'' remains open.  A non-Boolean OML is
not inconsistent in the proof-theoretic sense (no $\bot$ derivable),
but it is incompatible with single-context classical embedding.

\section*{The center obstruction (2026-05-16)}

Pt\'{a}k's construction reveals a striking fact: the bottom-right
obstruction can be \emph{exactly} the bottom-left obstruction (CE
failure on the Boolean center $C(L)$), just embedded in a non-Boolean
wrapper.  This suggests the admissibility condition for non-Boolean OMLs
might simply be CE applied to the center.

For $L(H)$: the center is trivial ($\{0, H\}$), so the center
obstruction is vacuous.  Gleason then provides $\sigma$-additivity
from lattice geometry (the continuum structure of projective space).

Tentative picture:
\begin{itemize}[nosep]
  \item Center non-trivial $\to$ CE on center is the obstruction
    (same mechanism as Boolean case)
  \item Center trivial (factor) $\to$ lattice geometry forces
    $\sigma$-additivity (Gleason)
  \item Center trivial but lattice too discrete $\to$ ??? (possible
    third obstruction, genuinely non-Boolean)
\end{itemize}

\section*{Questions for reading}

When reading Caramello, Vickers, Kalmbach, Hodges, Pt\'{a}k--Pulmannov\'{a}:
\begin{enumerate}[nosep]
  \item Does the author's formalism give a definition of coherence?
  \item Does the author distinguish consistency from coherence?
  \item Where would the author place CE in their framework?
  \item Does the non-Boolean case (quantum) fit naturally, or require
    a separate treatment?
  \item Is there a formal sense in which ``incompatibility'' (non-Boolean)
    is a species of ``inconsistency,'' or are they genuinely independent?
\end{enumerate}

\section*{Devil's advocate warnings (2026-05-16)}

\begin{enumerate}[nosep]
  \item \textbf{Q3 may be a category error.}  CE adds a hypothesis
    (to get $\sigma$-additivity).  Gleason has no hypothesis --- the
    lattice alone forces it.  A ``unified condition $C$'' that is
    substantive for Boolean algebras and vacuous for $L(H)$ is not
    a unification.  For Q3 to survive, one needs a mechanism-level
    (not just output-level) connection.
  \item \textbf{The square organizes, doesn't predict.}  Until it
    produces a theorem of the form ``given lattice property $X$,
    forcing holds iff $Y$,'' it is a filing system.
  \item \textbf{``Irreducibly algebraic'' is premature.}  One failed
    locale approach (Simpson) does not prove non-expressibility.
    Weaken to: tentative conclusion pending further investigation.
  \item \textbf{The Simpson conjecture is dead.}  $\delta_U$
    counterexample kills $\Leftarrow$; compact Hausdorff + Radon
    makes $\Rightarrow$ trivially true for all measures.  Locale
    non-spatiality is a constructive phenomenon that vanishes
    classically.
\end{enumerate}

\section*{Status}

Intuition pump.  Do not promote to stable files until a definition
crystallizes.  Hold lightly while reading.

The square is useful as a \emph{lens for reading} (organizes what to
look for in Kalmbach, Pt\'{a}k--Pulmannov\'{a}), not as a theorem
target.  The practical path remains: ship Papers~I and~II, close
Lean sorrys, consider formalization venues (ITP/CPP).

\end{document}
